% CS Labs Report --- Laboratory Work No. 2
\documentclass[12pt,a4paper]{article}

\usepackage[utf8]{inputenc}
\usepackage[T1]{fontenc}
\usepackage[english]{babel}
\usepackage{lmodern}
\usepackage[left=2cm,right=1.5cm,top=2cm,bottom=2cm]{geometry}
\usepackage{setspace}
\usepackage{graphicx}
\usepackage{xcolor}
\usepackage{listings}
\usepackage{titlesec}
\usepackage{caption}
\usepackage{float}
\usepackage{amsmath}
\usepackage{amssymb}
\usepackage{fancyhdr}
\usepackage{needspace}
\usepackage[hidelinks]{hyperref}

\graphicspath{{figures/}}
\onehalfspacing
\setlength{\parindent}{1.25cm}
\setlength{\parskip}{0pt}
\setlength{\emergencystretch}{2em}

\definecolor{codebackground}{HTML}{F3F4F6}
\definecolor{codeframe}{HTML}{D1D5DB}
\definecolor{commentgreen}{HTML}{3F6212}
\definecolor{placeholdertext}{HTML}{64748B}

\captionsetup[figure]{name=Fig,labelsep=period,font=small}

\titleformat{\section}
  {\normalfont\bfseries\centering\fontsize{13}{16}\selectfont}
  {\thesection}{0.8em}{\MakeUppercase}
\titleformat{name=\section,numberless}
  {\normalfont\bfseries\centering\fontsize{13}{16}\selectfont}
  {}{0pt}{\MakeUppercase}
\titleformat{\subsection}
  {\normalfont\bfseries\fontsize{12}{15}\selectfont}
  {\thesubsection}{0.7em}{}
\titlespacing*{\section}{0pt}{0pt}{12pt}
\titlespacing*{\subsection}{0pt}{10pt}{6pt}

\pagestyle{fancy}
\fancyhf{}
\fancyhead[L]{\small\textit{Cryptography and Security}}
\fancyhead[R]{\small Laboratory Work No. 2}
\fancyfoot[C]{\thepage}
\renewcommand{\headrulewidth}{0.4pt}
\renewcommand{\footrulewidth}{0pt}
\setlength{\headheight}{15pt}

\lstdefinestyle{go}{
  basicstyle=\ttfamily\small,
  backgroundcolor=\color{codebackground},
  frame=single,
  rulecolor=\color{codeframe},
  breaklines=true,
  columns=fullflexible,
  keepspaces=true,
  showstringspaces=false,
  commentstyle=\color{commentgreen}\itshape,
  morecomment=[l]{//},
  xleftmargin=0.4cm,
  xrightmargin=0.2cm,
  framexleftmargin=0.2cm,
  framesep=5pt,
  aboveskip=8pt,
  belowskip=8pt,
  extendedchars=true,
  literate=
    {Ă}{\u{A}}1 {Â}{\^{A}}1 {Î}{\^{I}}1
    {Ș}{\c{S}}1 {Ț}{\c{T}}1 {Ş}{\c{S}}1 {Ţ}{\c{T}}1
    {ă}{\u{a}}1 {â}{\^{a}}1 {î}{\^{i}}1
    {ș}{\c{s}}1 {ț}{\c{t}}1 {ş}{\c{s}}1 {ţ}{\c{t}}1
    {—}{\textemdash}1
}

\newcommand{\reportfigure}[4][\textwidth]{%
  \begin{figure}[H]
    \centering
    \includegraphics[width=#1]{#2}
    \caption{#3}
    \label{#4}
  \end{figure}
}

\begin{document}

\begin{titlepage}
  \thispagestyle{empty}
  \begin{center}
    \IfFileExists{figures/LOGO_UTM.jpg}{%
      \includegraphics[width=3.4cm]{figures/LOGO_UTM.jpg}\par
      \vspace{0.25cm}
    }{}

    {\bfseries\fontsize{14}{17}\selectfont
      MINISTRY OF EDUCATION, CULTURE AND RESEARCH\par
      OF THE REPUBLIC OF MOLDOVA\par
      Technical University of Moldova\par
      Faculty of Computers, Informatics and Microelectronics\par
      Department of Computer Science and Information Technologies\par
    }

    \vspace{1.4cm}
    {\bfseries\fontsize{16}{19}\selectfont Dinu Plesu, FAF-241\par}

    \vspace{1.4cm}
    {\bfseries\fontsize{36}{40}\selectfont Report\par}
    \vspace{0.25cm}
    {\itshape\fontsize{20}{24}\selectfont Laboratory Work No. 2\par}
    {\bfseries\itshape\fontsize{20}{24}\selectfont Cryptography and Security\par}

    \vspace{1cm}
    {\fontsize{15}{19}\selectfont Topic: Cryptanalysis of monoalphabetic ciphers --- frequency analysis attack (variant V20)\par}
  \end{center}

  \vfill
  \begin{flushright}
    {\fontsize{15}{18}\selectfont Checked by:}\par
    \textbf{Maia Zaica}
  \end{flushright}

  \vfill
  \begin{center}
    {\bfseries\fontsize{14}{17}\selectfont Chi\c{s}in\u{a}u --- 2026}
  \end{center}
\end{titlepage}

\pagenumbering{arabic}
\setcounter{page}{2}
\tableofcontents
\newpage

\section{Task}

A message intercepted from an enemy channel is known to have been encrypted with a monoalphabetic substitution cipher over a natural language. The variant assigned to the student ($k=20$) is the ciphertext V20. The task is to recover the plaintext without knowing the key, using the statistical attack described in the handout: build the letter-frequency table, apply the mechanical (most-frequent $\rightarrow$ most-frequent) substitution, then refine the substitution interactively by examining word patterns, short words and frequent digraphs until the text becomes readable.

\section{Theoretical contents}

In a monoalphabetic substitution cipher each letter of the alphabet is replaced by a fixed, unique substitute; the cipher is therefore a permutation of the alphabet. Although the key space is huge ($26!\approx4\times10^{26}$ for English), the frequency of each cipher letter is exactly the frequency of the corresponding plain letter. Every language has a characteristic frequency profile (Tables~2.1 and 2.2 of the handout for Romanian and English).

The attack proceeds iteratively. First the observed cipher frequencies are ranked and matched against the most frequent letters of the language --- a mechanical mapping that usually produces a partially readable text. The result is then refined by a human analyst: two- and three-letter words, recurring digraphs and the shapes of frequent words reveal further letter correspondences, which are applied and the text is re-examined. The quality of a candidate substitution can be measured with the chi-square statistic:

\begin{equation}
  \chi^2=\sum_i\frac{(O_i-E_i)^2}{E_i},
\end{equation}
where $O_i$ is the observed count of the $i$-th plain letter and $E_i$ is the expected count from the language frequency profile.

A substitution that approaches the true key drives $\chi^2$ down sharply and turns the masked word list (unknown letters shown as \texttt{*}) into recognisable words.

Nothing in the attack inherently requires a computer: it is a manual, human-driven procedure that can be carried out on paper from the frequency tables of the handout, or with any online frequency-analysis tool. A program only automates the bookkeeping --- counting letters, ranking them and keeping the candidate substitution consistent --- while every actual letter correspondence is still chosen by the analyst from the linguistic evidence.

\section{Implementation}

For reproducibility the attack was run with a small local tool, \texttt{lab2} (Go, file \texttt{freq.go}) --- the same manual procedure, with the arithmetic automated. The result is equally achievable by hand or with an online frequency-analysis website; the program was preferred only because it recomputes the frequency table and $\chi^2$ after every correspondence, keeps the substitution a bijection, and makes the whole attack repeatable from the repository. Option~2 asks for the language profile, prints the ranked frequency table of the ciphertext, applies the automatic frequency mapping, and enters an interactive loop with the commands \texttt{<CIFRATA>=<CLARA>} (assign a correspondence), \texttt{del}, \texttt{tab} (frequency table), \texttt{text} (current decryption) and \texttt{gata} (finish) --- the analyst, not the program, decides each mapping. Key fragments:

\begin{lstlisting}[style=go,caption={Key fragments of the automatic mapping and interactive refinement (\texttt{freq.go})}]
// Automatic mapping: rank cipher letters by frequency and match them
// to the most frequent letters of the language profile.
for idx, l := range observed {
    assign[l] = profile.letters[idx]
}

// Interactive refinement (bijective): assigning a plain letter frees it
// from any previous cipher letter, keeping the substitution a permutation.
for oc, om := range assign {
    if om == m && oc != c { delete(assign, oc) }
}
assign[c] = m
\end{lstlisting}

Because the intercepted text V20 is written in English, the English profile (Table~2.2) was selected for the analysis; with the automatic mapping in place, 13 manual corrections were enough to complete the substitution.

\section{Results of the work}

\reportfigure{fig1_caesar_warmup.png}{Option 1 --- a Caesar cryptogram produced with $k=11$ (warm-up example).}{fig:caesar-warmup}

\reportfigure{fig2_frequency_auto.png}{Ranked frequency table of V20 and the mechanical (automatic) substitution; $\chi^2$ of the automatic mapping $=47.8$.}{fig:frequency-auto}

\reportfigure{fig3_iteration_gn.png}{Interactive step: the correction \texttt{G=N} (the letter N is freed from its old mapping) and the partial decryption that follows it.}{fig:iteration-gn}

\reportfigure{fig4_plaintext.png}{The fully recovered plaintext (fragment): the Black Chamber and the Washington Naval Conference. The complete text (618 words) is saved to \texttt{Lab2/v20\_plain.txt}.}{fig:plaintext}

The recovered message is an excerpt from J.~Kahn's \textit{The Codebreakers} describing the Black Chamber, Yardley and the Washington Naval Conference; the complete text (618 words) was saved to \texttt{Lab2/v20\_plain.txt} and is reused as the signed message in Laboratory Work No.~6.

Read together, the four figures trace the full attack: Fig.~\ref{fig:frequency-auto} establishes the mechanical starting point --- the ranked cipher frequencies matched to English, with $\chi^2=47.8$ for the pure most-frequent $\rightarrow$ most-frequent guess; Fig.~\ref{fig:iteration-gn} shows a single linguistic correction (\texttt{G=N}) and the partial text it immediately unlocks; and Fig.~\ref{fig:plaintext} is the end state, a coherent plaintext. The 13 corrections between Fig.~\ref{fig:frequency-auto} and Fig.~\ref{fig:plaintext} are the analyst's reasoning, not the program's.

\section{Conclusions}

The monoalphabetic substitution cipher was fully broken without any knowledge of the key. Letter-frequency ranking plus word pattern analysis recovered the entire plaintext of variant V20 with 13 manual corrections on top of the automatic mapping. This confirms that monoalphabetic ciphers, whatever their key space, are insecure against a statistically grounded cryptanalytic attack --- the classical frequency analysis known since Al-Kindi. The attack is fundamentally interactive and linguistic: the analyst reads the text, recognises word shapes and steers each correspondence, while $\chi^2$ merely measures how close the candidate already is --- the software, whether a local script or an online tool, is only the workbench.

\section*{References}
\addcontentsline{toc}{section}{References}

\begin{enumerate}
  \item Lucrare de laborator nr.~2 --- Criptanaliza cifrurilor monoalfabetice. UTM, FCIM --- handout (\texttt{resources/}).
  \item Tables~2.1--2.4 --- reference letter frequencies for Romanian and English.
  \item J.~Kahn. \textit{The Codebreakers} --- the story of secret writing (source of the V20 text).
\end{enumerate}

\section*{Source Code}
\addcontentsline{toc}{section}{Source Code}

The complete source code of all laboratory works (Go sources for Labs~1--3, Python sources for Labs~4--6, together with the ciphertext and plaintext files used in Lab~2 and Lab~6) is available in the GitHub repository:

\begin{center}
  \url{https://github.com/Prince-of-Nothing/CS-labs.git}
\end{center}

\end{document}
